\chapter{Spaces, Quotients and Duality}\label{linear:ch:duality}

\begin{center}\itshape
I have every confidence that the effort I have applied to the science reported upon here, which has occupied a
considerable span of my lifetime and demanded the most intense exertions of my powers, is not to be lost.
\upshape(Hermann Grassmann, foreword to the \emph{Ausdehnungslehre} of 1862)
\end{center}
\noindent\emph{Start with Volume~1's linear algebra. The quotient and dual constructions here support Appendices~\ref{linear:ch:eigen} and~\ref{linear:ch:inner}; Volume~10 also uses the double dual.}

Grassmann's 1844 \emph{Ausdehnungslehre} asked readers to think beyond the geometry of three-dimensional space. He was a schoolteacher in Stettin, trained in theology and philology rather than university mathematics. Working largely alone, he developed objects that could be added and scaled in any number of dimensions. His book contained independence, subspaces, dimension, exterior products and the identity now called Grassmann's formula.

The mathematical ideas were powerful, but the philosophical presentation lost readers. M\"obius struggled to review it; Kummer praised its content while finding its form unsuitable. Grassmann's 1862 revision was more recognisably mathematical and still sold poorly. His later Sanskrit scholarship, including a Rig Veda translation and dictionary, received recognition sooner. Peano's 1888 account helped turn Grassmann's mathematical ideas into the vector-space axioms now taught at the start of linear algebra.

The useful objects here are spaces and maps, with matrices serving as coordinates. A quotient identifies vectors whose difference we agree to ignore. A dual space collects all linear measurements, and a transpose pulls measurements back along a map. These constructions also give a counting method: represent objects by independent vectors, then bound their number by the dimension.

\begin{think}
A town of $32$ residents permits clubs of odd size, provided distinct clubs overlap in an even number of residents. Can there be $32$ clubs? Could there be $33$? Encode membership by zeros and ones and try taking dot products modulo two. Section~\ref{linear:sec:dimarg} gives the bound.
\end{think}

\section{Spaces and bases}

Let $\mathbb K$ be a field, such as $\Q,\R,\C$ or $\mathbb F_p$. A \emph{vector space} is an abelian group $V$ with scalar multiplication satisfying
\[1v=v,\quad\lambda(\mu v)=(\lambda\mu)v,\quad
(\lambda+\mu)v=\lambda v+\mu v,\quad\lambda(v+w)=\lambda v+\lambda w.
\]
Examples include coordinate vectors, rectangular matrices, polynomials, bounded-degree polynomials and all functions $S\to\mathbb K$. The axioms make linear combinations meaningful in each of these settings.

A family is \emph{linearly independent}, or \emph{free}, when every finite vanishing linear combination has all coefficients zero. It \emph{spans} $V$ when finite combinations give every vector. A \emph{basis} does both, so each vector has exactly one finite coordinate expansion. The exchange lemma limits how many independent coordinates a finite spanning family can support.

\begin{lemma}[Steinitz exchange]\label{linear:lem:steinitz}
Suppose $v_1,\ldots,v_m$ are independent and belong to the span of $w_1,\ldots,w_n$. Then $m\le n$. After reordering the $w_j$, the family $v_1,\ldots,v_m,w_{m+1},\ldots,w_n$ spans the same space.
\end{lemma}
\begin{proof}
Induct on $m$, starting with the unchanged spanning family at $m=0$. After $m-1$ replacements, write
\[v_m=\sum_{i<m}\alpha_iv_i+\sum_{j\ge m}\beta_jw_j.
\]
Some $\beta_j$ must be nonzero; otherwise $v_m$ depends on its predecessors. Thus at least one $w_j$ remains, giving $m\le n$. Reorder to make $\beta_m\ne0$ and solve for $w_m$. Replacing it by $v_m$ leaves the span unchanged, completing the induction.
\end{proof}

\begin{theorem}[dimension]\label{linear:thm:dim}
A finitely spanned space $V$ has a finite basis, and all its bases have the same size, called $\dim V$. Every independent family extends to a basis, every spanning family contains a basis, and a family of $n$ vectors in an $n$-dimensional space is independent exactly when it spans.
\end{theorem}
\begin{proof}
Start with a finite spanning list and remove any vector dependent on the others until no removal is possible. The remaining list is a basis. Applying exchange in both directions makes any two finite bases equally large. No infinite independent family can occur: a finite spanning list of size $N$ forbids even an independent subfamily of size $N+1$.

Exchange now extends any independent family to a spanning list of the basis size. A spanning list of that size is independent, since removing a dependency would produce a smaller spanning list and hence a smaller basis. This proves extension and both directions of the final assertion. For an arbitrary spanning family, first extract finitely many elements spanning a fixed finite basis, then remove dependencies as above.
\end{proof}

For finite-dimensional subspaces $U,W$, begin with a basis of $U\cap W$ and extend it separately in $U$ and $W$. The intersection basis together with both sets of added vectors spans $U+W$. It is independent: in a vanishing combination, the added $U$ part lies in $U\cap W$, forcing its coefficients to vanish, and then independence in $W$ finishes. Counting gives
\[\dim(U+W)+\dim(U\cap W)=\dim U+\dim W.
\]

\begin{remark}
An infinite-dimensional space has no finite spanning family. Zorn's lemma (Volume~7, Chapter~1) still provides a basis, usually without an explicit description. The real numbers cannot have finite dimension over $\Q$, because $\Q^n$ is countable and $\R$ is not. Main Chapter~\ref{ch:banach} develops the topological issues that arise in infinite dimension; most arguments in this appendix remain finite-dimensional.
\end{remark}

\section{Sums, complements and quotients}\label{linear:sec:quotients}

Subspaces form a \emph{direct sum} when each vector in their sum has a unique decomposition into one vector from each. For two subspaces this is equivalent to intersection $\{0\}$. A \emph{complement} $C$ of $W$ in $V$ satisfies $V=W\oplus C$. In finite dimension, extend a basis of $W$ and span the added vectors to find one. Complements are usually not unique: any other line complements a fixed line in a plane. A quotient records the remaining directions without choosing a complement.

\begin{definition}\label{linear:def:quotient}
For a subspace $W\subseteq V$, write $v\equiv v'\pmod W$ when $v-v'\in W$. Its equivalence classes are the affine sets $v+W$. The \emph{quotient space} $V/W$ consists of these classes with operations
\[(v+W)+(v'+W)=(v+v')+W,\qquad\lambda(v+W)=\lambda v+W.\]
\end{definition}

Changing either representative changes these expressions only by an element of $W$, so the operations are well defined. The vector-space axioms descend from $V$. The projection $\pi(v)=v+W$ is linear, surjective and has kernel $W$. In Figure~\ref{linear:fig:quotient}, parallel lines are individual quotient vectors; a complement selects one representative on each.
\begin{figure}[ht]
\centering
\fig{linear-ch01-quotient}
\caption{The quotient of the plane by a line $W$: its elements are the lines parallel to $W$. Any complement $C$ meets
each of them once, so $C\to V/W$ is an isomorphism; the quotient itself does not depend on the choice of $C$.}
\label{linear:fig:quotient}
\end{figure}

\begin{theorem}[quotients]\label{linear:thm:quotient}
Let $W$ be a subspace of finite-dimensional $V$.
\begin{enumerate}[label=(\alph*)]
\item Every complement $C$ maps isomorphically onto $V/W$ under $\pi|_C$, and $\dim(V/W)=\dim V-\dim W$.
\item If a linear $f:V\to U$ vanishes on $W$, there is exactly one linear $\bar f:V/W\to U$ with $f=\bar f\circ\pi$.
\item Every linear map induces an isomorphism $V/\ker f\cong\operatorname{Im}f$; consequently $\dim V=\dim\ker f+\operatorname{rank}f$.
\item For $W\subseteq W'\subseteq V$, the image $W'/W$ is a subspace of $V/W$ and $(V/W)/(W'/W)\cong V/W'$.
\end{enumerate}
\end{theorem}
\begin{proof}
(a) A vector in the kernel of $\pi|_C$ belongs to $C\cap W=\{0\}$. Every class has a representative in $C$, since $v=w+c$. This proves the isomorphism and its dimension formula.

(b) The equation forces $\bar f(v+W)=f(v)$. It is well defined because $v-v'\in W$ implies $f(v)=f(v')$; linearity descends from $f$. Surjectivity of $\pi$ makes the factor unique.

(c) Apply (b) with $W=\ker f$ and target $\operatorname{Im}f$. The induced map is onto, and a zero image means $v\in\ker f$, the zero quotient class. Part (a) gives rank--nullity.

(d) The map $v+W\mapsto v+W'$ is well defined, linear and onto. Its kernel is precisely $W'/W$. Part (c) gives the asserted quotient isomorphism.
\end{proof}

The factorisation and first-isomorphism arguments themselves use no finite-dimensional assumption; that assumption is needed for the dimension formulas. This lets us use them on polynomial and function spaces as well.

\begin{example}\label{linear:ex:quotients}
(1) For a nonzero polynomial $P$ of degree $d$, Euclidean division gives exactly one representative of degree less than $d$ in every class of $\mathbb K[X]/(P)$. Thus the quotient has dimension $d$, even though $\mathbb K[X]$ has infinite dimension. When $P$ is constant the quotient is zero, with the zero remainder understood separately. Over $\R$, the relation $X^2+1=0$ makes the class of $X$ square to $-1$, identifying this quotient with $\C$, as in Cauchy's 1847 construction.

(2) For $T\subseteq S$, restrict functions from $S$ to $T$. This map is onto, since a function on $T$ can be extended by zero. Its kernel consists of functions vanishing on $T$, so quotienting by that kernel gives $\mathbb K^T$.

(3) For $A\in M_{m,n}(\mathbb K)$, the system $Ax=b$ is solvable exactly when $b\in\operatorname{Im}A$. Equivalently, the class of $b$ vanishes in the \emph{cokernel} $\mathbb K^m/\operatorname{Im}A$. Its dimension $m-\operatorname{rank}A$ counts independent compatibility conditions, which Section~\ref{linear:sec:transpose} will identify explicitly.
\end{example}

\section{The dual space}
\begin{definition}\label{linear:def:dual}
The \emph{dual space} $V^*$ consists of all linear maps $V\to\mathbb K$, called \emph{linear forms}, with pointwise addition and scaling. For a basis $e_1,\ldots,e_n$, its coordinate forms are specified by $e_i^*(e_j)=\delta_{ij}$.
\end{definition}

Forms measure a vector by one scalar. Examples are a coordinate, a matrix trace, evaluation of a polynomial at a fixed point, or integration on a suitable function space. Linearity is what these operations share.

\begin{proposition}[dual basis]\label{linear:prop:dualbasis}
The coordinate forms $e_i^*$ form a basis of $V^*$, and every form has the expansion $\varphi=\sum_i\varphi(e_i)e_i^*$. In particular, $\dim V^*=\dim V$.
\end{proposition}
\begin{proof}
The expansion has the same values as $\varphi$ on every basis vector, hence on all vectors. This proves spanning. Evaluating a relation $\sum_i\lambda_ie_i^*=0$ at $e_j$ gives $\lambda_j=0$, proving independence.
\end{proof}

Matching a basis to its dual gives an isomorphism $V\cong V^*$, but changing the basis changes that matching. Over $\R$, replacing $e_1$ by $2e_1$ replaces its coordinate form by $e_1^*/2$. An inner product can supply an additional identification (Appendix~\ref{linear:ch:inner}). The double dual needs no such choice: a vector evaluates forms at itself.

\begin{theorem}[the double dual]\label{linear:thm:bidual}
For finite-dimensional $V$, the map $\iota:V\to V^{**}$ defined by $\iota(v)(\varphi)=\varphi(v)$ is an isomorphism independent of a basis choice.
\end{theorem}
\begin{proof}
Both the evaluation in $\varphi$ and the dependence on $v$ are linear. If $v\ne0$, extend it to a basis and take the first coordinate form, which evaluates to one on $v$. Thus $\iota$ is injective. Applying the dual-dimension proposition twice gives equal dimensions, so it is surjective.
\end{proof}

\begin{remark}
The algebraic evaluation map remains injective in infinite dimension but is no longer onto. To see the obstruction, choose an infinite basis indexed by $I$, so $V^*=\mathbb K^I$. By basis extension, a linear functional on $\mathbb K^I$ can vanish on all finitely supported sequences and take value one on the constant sequence one. It cannot be a finite combination of coordinate evaluations: testing the coordinate unit sequences would make every coefficient zero. It therefore lies outside $\iota(V)$. For $V=\Q[X]$, the coefficient space $V$ is countable whereas its dual is the uncountable sequence space $\Q^{\N}$. Main Chapter~\ref{ch:banach} instead uses continuous forms to obtain a topological dual theory.
\end{remark}

A dual basis can also make coordinates easier to find. Polynomial interpolation is a useful example.

\begin{example}[Lagrange interpolation]\label{linear:ex:lagrange}
Choose distinct $x_0,\ldots,x_n\in\mathbb K$ and let $V=\mathbb K[X]_{\le n}$. Evaluation forms $\varepsilon_i(P)=P(x_i)$ pair with
\[L_j(X)=\prod_{k\ne j}\frac{X-x_k}{x_j-x_k},\qquad\varepsilon_i(L_j)=\delta_{ij}.
\]
Evaluating any vanishing combination proves the $L_j$ independent. There are $n+1$ of them, so they form a basis with dual basis $\varepsilon_i$. Thus $P=\sum_jP(x_j)L_j$, and arbitrary prescribed values at these nodes have exactly one interpolant of degree at most $n$.
\end{example}

For rectangular matrices, write $\tau_A(X)=\operatorname{tr}(AX)$ with $A\in M_{n,m}(\mathbb K)$ and $X\in M_{m,n}(\mathbb K)$.
\begin{proposition}[the trace pairing]\label{linear:prop:trace}
The map $A\mapsto\tau_A$ identifies $M_{n,m}(\mathbb K)$ isomorphically with $M_{m,n}(\mathbb K)^*$. Thus every form on this matrix space has a unique trace representation.
\end{proposition}
\begin{proof}
Since $\operatorname{tr}(AX)=\sum_{i,j}a_{ij}x_{ji}$, evaluating at $E_{ji}$ recovers $a_{ij}$. The map is injective and linear between spaces of dimension $mn$, hence bijective.
\end{proof}

\section{Annihilators and the transpose}\label{linear:sec:transpose}
\begin{definition}
For $S\subseteq V$, its \emph{annihilator} is $S^\circ=\{\varphi\in V^*:\varphi(s)=0\text{ for every }s\in S\}$. For a linear map $f:V\to U$, its \emph{transpose} is $f^{\mathsf T}:U^*\to V^*$, $\psi\mapsto\psi\circ f$.
\end{definition}

\begin{theorem}[annihilators]\label{linear:thm:annihilator}
Let $V$ have finite dimension $n$ and $W\subseteq V$ be a subspace.
\begin{enumerate}[label=(\alph*)]
\item There is a canonical isomorphism $W^\circ\cong(V/W)^*$, and $\dim W^\circ=n-\dim W$.
\item Identifying $V$ with $V^{**}$, one has $W^{\circ\circ}=W$. In particular, a subspace of codimension $k$ is the common kernel of $k$ forms, an intersection of $k$ hyperplanes.
\item Annihilators exchange sums and intersections:
\[(W_1+W_2)^\circ=W_1^\circ\cap W_2^\circ,\qquad
(W_1\cap W_2)^\circ=W_1^\circ+W_2^\circ.\]
\end{enumerate}
\end{theorem}
\begin{proof}
(a) Pull a form on $V/W$ back by $\pi$. This gives exactly the forms vanishing on $W$, with inverse given by factorisation through the quotient. The quotient and dual dimensions give the count.

(b) Every $w\in W$ annihilates the forms in $W^\circ$, so $W\subseteq W^{\circ\circ}$. Applying (a) twice makes their dimensions equal. Taking a basis of $W^\circ$ then expresses $W$ as the common kernel of those $k$ forms.

(c) A form kills a sum precisely when it kills both summands, proving the first identity. The sum of the two annihilators lies in the annihilator of the intersection. Grassmann's formula, (a) and the first identity show its dimension is
\[(n-\dim W_1)+(n-\dim W_2)-(n-\dim(W_1+W_2))
 =n-\dim(W_1\cap W_2).
\]
Equal dimensions prove the second identity.
\end{proof}

\begin{theorem}[the transpose]\label{linear:thm:transpose}
For a linear map $f:V\to U$ between finite-dimensional spaces:
\begin{enumerate}[label=(\alph*)]
\item its matrix $A$ in chosen bases becomes $A^{\mathsf T}$ in the corresponding dual bases;
\item $\ker f^{\mathsf T}=(\operatorname{Im}f)^\circ$ and $\operatorname{Im}f^{\mathsf T}=(\ker f)^\circ$;
\item $\operatorname{rank}f^{\mathsf T}=\operatorname{rank}f$, so a matrix has equal row and column ranks.
\end{enumerate}
\end{theorem}
\begin{proof}
(a) If $f(e_j)=\sum_i a_{ij}u_i$, the coefficient in row $j$, column $i$ of the transpose map is $f^{\mathsf T}(u_i^*)(e_j)=u_i^*(f(e_j))=a_{ij}$.

(b) The equation $\psi\circ f=0$ says exactly that $\psi$ kills $\operatorname{Im}f$, giving the kernel identity. Every pullback form kills $\ker f$, giving one image inclusion. The kernel identity and the annihilator dimension yield $\operatorname{rank}f^{\mathsf T}=\operatorname{rank}f$. This equals $\dim(\ker f)^\circ$, so the inclusion is equality. The rank computation also proves (c).
\end{proof}

Row rank equals column rank because pulling measurements back retains the rank of the original map. The kernel identity gives the exact compatibility test for a linear system.

\begin{corollary}[Fredholm alternative]\label{linear:cor:fredholm}
For $A\in M_{m,n}(\mathbb K)$ and $b\in\mathbb K^m$, the equation $Ax=b$ is solvable exactly when $y^{\mathsf T}b=0$ for every $y$ satisfying $A^{\mathsf T}y=0$. Exactly one alternative holds: a solution $x$ exists, or a vector $y$ satisfies $A^{\mathsf T}y=0$ and $y^{\mathsf T}b=1$.
\end{corollary}
\begin{proof}
Identify forms on $\mathbb K^m$ with row vectors. The transpose theorem gives $\ker A^{\mathsf T}=(\operatorname{Im}A)^\circ$, whose annihilator is $\operatorname{Im}A$. This proves the test. A failed test gives $y^{\mathsf T}b\ne0$; rescale $y$ to make it one. Such a certificate cannot coexist with a solution, since then $y^{\mathsf T}b=y^{\mathsf T}Ax=0$.
\end{proof}

A solution certifies consistency by substitution. The other certificate combines the equations into $0=1$, certifying inconsistency. This finite-dimensional duality also underlies linear programming and Delsarte's bounds in Volume~3.
\begin{example}[Lights Out]\label{linear:ex:lightsout}
Pressing a cell of a Lights Out board toggles that cell and its horizontal and vertical neighbours. Over $\mathbb F_2$, a press pattern $x$ changes the board by $Ax$, where $A$ has ones for self and neighbour incidences. Order is irrelevant and double presses cancel. Thus switching off $b$ is the linear problem $Ax=b$, with symmetric $A$.

On the $5\times5$ board the kernel has dimension two, with independent quiet patterns $q_1,q_2$. A board $b$ is solvable exactly when $q_1\cdot b=q_2\cdot b=0$, so one quarter of configurations are solvable and each has four press patterns. The $3\times3$ and $6\times6$ boards are invertible; the $4\times4$ board has kernel dimension four. The row argument below proves these claims exactly.
\end{example}

Here is an exact way to obtain the board ranks. On an $n\times n$ board, let $B_n$ have ones on its diagonal and adjacent diagonals. A quiet press pattern with row vectors $r_j$ satisfies
\[r_{j+1}=B_nr_j+r_{j-1},\qquad r_0=r_{n+1}=0.
\]
Hence it is determined by its first row. Define matrix polynomials $P_0=0,P_1=I$, $P_{j+1}=B_nP_j+P_{j-1}$. The allowable first rows are exactly $\ker P_{n+1}$. For $n=5$,
\[P_6=B_5^5+B_5=
\begin{pmatrix}0&1&1&0&1\\1&1&1&0&0\\1&1&0&1&1\\0&0&1&1&1\\1&0&1&1&0\end{pmatrix}.
\]
Its kernel consists of $(b+c,b,c,b,b+c)$, as its rows show directly. Propagating the choices $(b,c)=(1,0),(0,1)$ gives two independent quiet patterns $q_1,q_2$ (Figure~\ref{linear:fig:lightsout} shows one). Therefore $\dim\ker A=2$, and symmetry plus Fredholm gives
\[b\text{ solvable}\iff q_1\cdot b=q_2\cdot b=0\quad\text{in }\mathbb F_2.
\]
The two independent checks leave $2^{23}$ boards, one quarter of all configurations. Each solution set is a coset of the two-dimensional kernel and has four press patterns.

The same recurrence gives $P_4=B_3^3$ with determinant one, $P_5=B_4^4+B_4^2+I=0$, and $P_7=B_6^6+B_6^4+I$ with determinant one, all in $\mathbb F_2$. These are finite matrix identities obtained by the stated recurrence and elimination. For an explicit check of the last determinant, its rows are
\[(1,0,0,0,1,0),\ (0,1,0,1,0,1),\ (0,0,0,0,1,0),\]
\[(0,1,0,0,0,0),\ (1,0,1,0,1,0),\ (0,1,0,0,0,1).
\]
Their homogeneous equations successively force coordinates $5,2,1,6,4,3$ to zero. Thus the $3\times3$ and $6\times6$ boards are invertible, while the $4\times4$ board has kernel dimension four. Exercise~\ref{linear:ex:lights3} gives explicit single-light solutions for the smallest case. The electronic game was marketed by Tiger Electronics in 1995, but the parity algebra does not depend on the device.
\begin{figure}[ht]
\centering
\fig{linear-ch01-lightsout}
\caption{Lights Out. Left: pressing the marked button toggles it and its four neighbours. Right: one of the two quiet
patterns of the $5\times5$ board; pressing all its lit buttons changes nothing, and every solvable configuration has an
even number of lit lights in common with it.}
\label{linear:fig:lightsout}
\end{figure}

\section{The dimension argument}\label{linear:sec:dimarg}

To bound a family's size, represent its members by independent vectors in a finite-dimensional space. The resulting upper bound is simply the dimension. Volume~3 uses this method throughout combinatorics; here are three examples.

\begin{theorem}[Oddtown; Berlekamp, 1969]\label{linear:thm:oddtown}
Suppose subsets $C_1,\ldots,C_m$ of $\{1,\ldots,n\}$ have odd sizes and distinct pairs have even intersection sizes. Then $m\le n$.
\end{theorem}
\begin{proof}
Their indicators $v_i\in\mathbb F_2^n$ satisfy $v_i\cdot v_j=\delta_{ij}$. Taking the dot product of $\sum_i\lambda_iv_i=0$ with $v_j$ yields $\lambda_j=0$. They are independent, so at most $n$ can occur.
\end{proof}

The $32$ singleton clubs meet the town's rules, and the theorem excludes $33$. Evaluation made the Lagrange polynomials independent in the same way: measurements gave an identity pairing matrix. A triangular pairing matrix also suffices.

\begin{proposition}[triangular criterion]\label{linear:prop:triangular}
Let $v_1,\ldots,v_m\in V$ and $\varphi_1,\ldots,\varphi_m\in V^*$ satisfy $\varphi_i(v_i)\ne0$ and $\varphi_i(v_j)=0$ for $j<i$. Then both listed families are independent.
\end{proposition}
\begin{proof}
In a nontrivial vector relation, take the largest index $k$ with nonzero coefficient and apply $\varphi_k$. All other terms vanish, a contradiction. In a nontrivial relation between forms, take the smallest nonzero coefficient index $k$ and evaluate at $v_k$. The terms with larger indices vanish by triangularity, and those with smaller indices have zero coefficient. Again the diagonal term cannot vanish.
\end{proof}

A \emph{hyperplane} is a codimension-one subspace, equivalently the kernel of a nonzero form by Theorem~\ref{linear:thm:annihilator}. The trace pairing lets us recognise a matrix hyperplane without choosing its basis.

\begin{theorem}\label{linear:thm:hyperplane}
For every $n\ge2$, each hyperplane in $M_n(\mathbb K)$ contains an invertible matrix.
\end{theorem}
\begin{proof}
Write the hyperplane as $\{X:\operatorname{tr}(AX)=0\}$ with $A\ne0$. Row and column elimination give $A=PJ_rQ$, where $P,Q$ are invertible and $J_r$ is diagonal with $r$ ones and the rest zero. Let $C$ cyclically permute the $n$ coordinate vectors. Since $n\ge2$, it has zero diagonal and is invertible. Then $X=Q^{-1}CP^{-1}$ is invertible and
\[\operatorname{tr}(AX)=\operatorname{tr}(PJ_rCP^{-1})=\operatorname{tr}(J_rC)=0.
\]
\end{proof}

At $n=1$ the only hyperplane is $\{0\}$, so the restriction is necessary. Although singular matrices satisfy one polynomial equation, they cannot contain a whole linear hyperplane when $n\ge2$. Exercise~\ref{linear:ex:singularspace} asks for large singular subspaces instead.

\begin{proposition}[unisolvent points]\label{linear:prop:unisolvent}
For an $n$-dimensional space $V$ of functions $S\to\mathbb K$, there exist $s_1,\ldots,s_n\in S$ such that every prescribed value list $y_1,\ldots,y_n$ has exactly one interpolant in $V$.
\end{proposition}
\begin{proof}
The evaluations $\varepsilon_s$ span $V^*$. Indeed, their span has annihilator consisting of functions zero at every point, hence only the zero function; the annihilator dimension theorem makes the span all of $V^*$. Extract an evaluation basis $\varepsilon_{s_1},\ldots,\varepsilon_{s_n}$. The map $f\mapsto(f(s_1),\ldots,f(s_n))$ is injective, since its zero vectors are killed by every form. Equal dimensions make it bijective. For $n=0$, the empty value list has the unique zero interpolant.
\end{proof}

Any $n+1$ distinct nodes work for univariate polynomials of degree at most $n$. Likewise any $2N+1$ distinct circle points determine a trigonometric polynomial of degree $N$ (Exercise~\ref{linear:ex:trig}). In two variables, six points on a circle fail for degree-two polynomials, since the circle's nonzero quadratic equation vanishes there. Some other choice of six points must work; an added exercise constructs one.

\begin{beauty}
Annihilators reverse inclusion and exchange sums with intersections. The double annihilator returns the original finite-dimensional subspace. Pulling forms back turns a map around and preserves its rank. These are concrete dualities: the same structure can be studied through vectors or through measurements. Projective geometry in Volume~9 and optimisation provide further versions with their own hypotheses.
\end{beauty}

\begin{figure}[ht]
\centering\fig{appA-rw-evaluation}
\caption{On cubic polynomials, evaluation at zero and one has kernel spanned by $X^2-X$ and $X^3-X$. The quotient keeps exactly the two endpoint values; the linear interpolant selects a representative.}
\label{fig:v8appAeval}
\end{figure}
\begin{figure}[ht]
\centering\fig{appA-rw-certificate}
\caption{The map $(x_1,x_2)\mapsto(x_1,x_2,x_1+x_2)$ has image annihilated by $(-1,-1,1)$. The scalar $b_3-b_1-b_2$ is the compatibility test for its three equations.}
\label{fig:v8appAcert}
\end{figure}
\section{Exercises}
\begin{exercise}[$\star$]\label{linear:ex:logs}
Prove rational independence of $\log2,\log3,\log5$, then of the logarithms of all primes. Infer another proof that $\R$ has infinite dimension over $\Q$.
\end{exercise}
\begin{exercise}[$\star$]\label{linear:ex:grassmann}
For two six-dimensional subspaces of a ten-dimensional space, determine every possible intersection dimension. Realise $U+W$ as $(U\oplus W)/D$ with $D\cong U\cap W$, and recover Grassmann's identity by rank--nullity.
\end{exercise}
\begin{exercise}[$\star$]\label{linear:ex:quotientdual}
For a subspace $W$ of finite-dimensional $V$, restrict forms to construct the canonical isomorphism $V^*/W^\circ\cong W^*$.
\end{exercise}
\begin{exercise}[$\star\star$]\label{linear:ex:traceform}
(a) Prove that a form on $M_n(\mathbb K)$, $n\ge1$, satisfies $\varphi(XY)=\varphi(YX)$ for all matrices exactly when it is a multiple of trace. (b) Show that commutators $XY-YX$ span the trace-zero hyperplane.
\end{exercise}
\begin{exercise}[$\star$]\label{linear:ex:symanti}
Under $\operatorname{tr}(AX)$ on real square matrices, identify the annihilators of the symmetric subspace and the upper triangular subspace.
\end{exercise}
\begin{exercise}[$\star\star$]\label{linear:ex:shift}
Classify the forms on $\R[X]_{\le n}$ unchanged by $P(X)\mapsto P(X+1)$, for every $n\ge0$.
\end{exercise}
\begin{exercise}[$\star\star$]\label{linear:ex:quadrature}
(a) At distinct nodes $x_0,\ldots,x_n\in[0,1]$, prove existence and uniqueness of weights integrating every polynomial of degree at most $n$ exactly. (b) Compute the weights at $0,1/2,1$ and verify exactness through degree three: Simpson's rule.
\end{exercise}
\begin{exercise}[$\star$]\label{linear:ex:trig}
For complex trigonometric polynomials $\sum_{|k|\le N}c_ke^{ikx}$, prove unique interpolation at any $2N+1$ distinct points of $[0,2\pi)$.
\end{exercise}
\begin{exercise}[$\star\star$]\label{linear:ex:lights3}
On a $3\times3$ Lights Out board, construct press patterns toggling each individual light alone. Use only exact parity arguments, then explain why this solves every configuration.
\end{exercise}
\begin{exercise}[$\star\star$]\label{linear:ex:eventown}
\emph{Reverse parity rules.} Suppose $C_1,\ldots,C_m\subseteq\{1,\ldots,n\}$ have even sizes and odd pairwise intersections. Prove $m\le n+1$, improving to $m\le n$ for even $m$, by studying $J-I$ over $\mathbb F_2$.
\end{exercise}
\begin{exercise}[$\star\star$]\label{linear:ex:singularspace}
(a) Construct a singular-matrix subspace of dimension $n(n-1)$. (b) If every nonzero matrix in a subspace of $M_n(\mathbb K)$ is invertible, bound its dimension by $n$. Attain dimension two in $M_2(\R)$.
\end{exercise}
\begin{exercise}[$\star\star$]\label{linear:ex:hyperplanes}
For $\varphi,\varphi_1,\ldots,\varphi_k\in V^*$, prove that $\varphi$ belongs to their span exactly when $\bigcap_{i=1}^k\ker\varphi_i\subseteq\ker\varphi$. You may take $V$ finite-dimensional as in the annihilator theorem; determine why the finite family of forms also permits the result for arbitrary $V$.
\end{exercise}
\begin{exercise}[$\star\star\star$]\label{linear:ex:hamel}
Choose a basis of $\R$ over $\Q$ using Zorn (Volume~7, Chapter~1). Construct an additive function $f:\R\to\R$ that is not $x\mapsto cx$ for any real $c$. Prove it is unbounded on every nondegenerate interval. Single-point intervals are excluded from this claim.
\end{exercise}
\begin{exercise}[$\star\star\star$]\label{linear:ex:fisher}
\emph{Constant intersections.} For $n\ge1$, let $C_1,\ldots,C_m$ be distinct subsets of $\{1,\ldots,n\}$ with $|C_i\cap C_j|=\lambda\ge1$ whenever $i\ne j$. Show $m\le n$ using a real Gram matrix.
\end{exercise}

\begin{exercise}[$\star\star$]\label{ex:v8cAa}
\emph{Cubic polynomials modulo endpoint zeros.} In $V=\R[X]_{\le3}$ let $W=\operatorname{span}(X^2-X,X^3-X)$. Find $W^\circ$ and identify $V/W$ with $\R^2$ by endpoint evaluations. Give the unique linear representative of each class and reduce $1+2X+3X^2+4X^3$.
\end{exercise}
\begin{exercise}[$\star\star$]\label{ex:v8cAb}
\emph{All complements as graphs.} Fix $V=W\oplus C$. Prove that every complement of $W$ is uniquely $\{L(c)+c:c\in C\}$ for a linear $L:C\to W$. Over $\mathbb F_q$, deduce that a $k$-dimensional subspace of an $n$-dimensional space has $q^{k(n-k)}$ complements.
\end{exercise}
\begin{exercise}[$\star$]\label{ex:v8cAc}
\emph{Three equations and one certificate.} For $A(x_1,x_2)=(x_1,x_2,x_1+x_2)$ over any field, find $\ker A^{\mathsf T}$. Classify solvable $Ax=b$ and give a certificate with $y^{\mathsf T}b=1$ whenever it is unsolvable.
\end{exercise}
\begin{exercise}[$\star\star$]\label{ex:v8cAd}
\emph{Six explicit interpolation points.} Over a field of characteristic different from two, prove that $(0,0),(1,0),(2,0),(0,1),(1,1),(0,2)$ are unisolvent for bivariate polynomials of total degree at most two. Explain where the characteristic assumption enters.
\end{exercise}

\section{Hints}
\begin{hint}{linear:ex:logs}
Clear rational denominators, exponentiate and compare prime exponents.
\end{hint}
\begin{hint}{linear:ex:grassmann}
The addition map from the external direct sum has kernel $\{(x,-x):x\in U\cap W\}$.
\end{hint}
\begin{hint}{linear:ex:quotientdual}
Extend a form by making it zero on a complement. The restriction kernel is its annihilator.
\end{hint}
\begin{hint}{linear:ex:traceform}
Use the trace representation and force its matrix to commute with every elementary matrix.
\end{hint}
\begin{hint}{linear:ex:symanti}
Test $E_{ii}$ and $E_{ij}+E_{ji}$ for the first subspace, then $E_{ij}$ with $i\le j$ for the second.
\end{hint}
\begin{hint}{linear:ex:shift}
The form must kill the image of $\Delta P=P(X+1)-P(X)$. Find that image, treating constants separately.
\end{hint}
\begin{hint}{linear:ex:quadrature}
Integrate each Lagrange basis polynomial to get its weight.
\end{hint}
\begin{hint}{linear:ex:trig}
Multiplication by $e^{iNx}$ converts the expression to a degree-$2N$ polynomial in $e^{ix}$.
\end{hint}
\begin{hint}{linear:ex:lights3}
Square symmetries reduce single lights to a corner, an edge centre and the board centre. Add their solutions modulo two.
\end{hint}
\begin{hint}{linear:ex:eventown}
The Gram rank cannot exceed $n$. Squaring $J-I$ distinguishes even and odd $m$.
\end{hint}
\begin{hint}{linear:ex:singularspace}
Force the first column to vanish. For the upper bound, evaluate every matrix in the subspace at $e_1$.
\end{hint}
\begin{hint}{linear:ex:hyperplanes}
Let $F(v)=(\varphi_1(v),\ldots,\varphi_k(v))$. A form killing $\ker F$ descends to its finite-dimensional image.
\end{hint}
\begin{hint}{linear:ex:hamel}
Prescribe incompatible scaling ratios at two basis vectors. Boundedness on an interval forces continuity of an additive function.
\end{hint}
\begin{hint}{linear:ex:fisher}
If $m\ge2$, the Gram matrix is $\lambda J+\operatorname{diag}(|C_i|-\lambda)$, with at most one zero diagonal correction. Handle $m\le1$ first.
\end{hint}
\begin{hint}{ex:v8cAa}
The two generators vanish at both endpoints. Rank--nullity identifies the full evaluation kernel.
\end{hint}
\begin{hint}{ex:v8cAb}
Projection to $C$ restricts to an isomorphism on any complement of $W$.
\end{hint}
\begin{hint}{ex:v8cAc}
A row vector kills the image precisely when its first two coordinates are the negatives of its third.
\end{hint}
\begin{hint}{ex:v8cAd}
First force the restriction to each coordinate axis to vanish; the remaining mixed term is detected at $(1,1)$.
\end{hint}

\section{Solutions}
\begin{solution}{linear:ex:logs}
\idea Rational relations between prime logarithms would give two different prime factorizations of one integer.
\writeup Clear denominators in a relation $\sum_iq_i\log p_i=0$, making every coefficient integral. Exponentiation gives $\prod_{q_i>0}p_i^{q_i}=\prod_{q_i<0}p_i^{-q_i}$. Unique factorisation forces every exponent to vanish. Infinitely many primes therefore give arbitrarily large independent subsets of $\R$ over $\Q$, excluding finite dimension by Steinitz.
\end{solution}
\begin{solution}{linear:ex:grassmann}
\idea Rank--nullity for addition counts the repeated intersection directions exactly once.
\writeup The map $(u,w)\mapsto u+w$ has image $U+W$ and kernel $D=\{(x,-x):x\in U\cap W\}$. The first isomorphism theorem gives the quotient description and Grassmann's formula. The intersection dimension $k$ lies between $12-10=2$ and $6$. Every such integer occurs: with a coordinate basis of the ten-dimensional space, take $U=\operatorname{span}(e_1,\ldots,e_6)$ and $W=\operatorname{span}(e_{7-k},\ldots,e_{12-k})$. Their common basis vectors are exactly $e_{7-k},\ldots,e_6$.
\end{solution}
\begin{solution}{linear:ex:quotientdual}
\idea Restriction loses exactly the forms that vanish on $W$.
\writeup Restriction $r:V^*\to W^*$ has kernel $W^\circ$. Every form on $W$ extends by zero on a chosen complement, so $r$ is onto. The first isomorphism theorem gives $V^*/W^\circ\cong W^*$. Restriction defines the isomorphism without a choice; the complement only proves surjectivity.
\end{solution}
\begin{solution}{linear:ex:traceform}
\idea Trace nondegeneracy converts a condition on forms into commutation with all matrices.
\writeup (a) Write $\varphi(X)=\operatorname{tr}(AX)$. Cyclicity of trace rewrites the hypothesis as $\operatorname{tr}((AX-XA)Y)=0$ for all $X,Y$. Nondegeneracy forces $AX=XA$ for all $X$. Commuting with $E_{ii}$ forces $A$ diagonal; commuting with $E_{ij}$ forces all its diagonal entries equal. For $n=1$ the latter condition is vacuous and the same scalar conclusion holds. Thus $\varphi=c\operatorname{tr}$, and every such form works.

(b) The annihilator of the commutator span consists exactly of these forms. Trace is nonzero, since it takes value one at $E_{11}$ even in positive characteristic. The annihilator has dimension one, so its double annihilator is precisely the trace-zero subspace of dimension $n^2-1$.
\end{solution}
\begin{solution}{linear:ex:symanti}
\idea Elementary basis tests expose the matrix entries allowed in each annihilator.
\writeup Annihilating $E_{ii}$ gives $a_{ii}=0$; annihilating $E_{ij}+E_{ji}$ gives $a_{ij}+a_{ji}=0$. Thus symmetric matrices have the skew-symmetric annihilator. For upper triangular matrices, $\operatorname{tr}(AE_{ij})=a_{ji}=0$ whenever $i\le j$. This forces all entries on and below the diagonal to vanish, leaving the strictly upper triangular matrices.
\end{solution}
\begin{solution}{linear:ex:shift}
\idea Translation invariance means annihilating the finite-difference image.
\writeup For $n\ge1$, $\Delta X^k=kX^{k-1}+\cdots$ lowers every nonconstant degree by one. Its kernel is the constants, so its image has dimension $n$ and equals $\R[X]_{\le n-1}$. A form killing that image is exactly $P\mapsto c[X^n]P$. When $n=0$, $\Delta=0$ and every form on the constant space has the same description $c[X^0]P$.
\end{solution}
\begin{solution}{linear:ex:quadrature}
\idea Evaluation is a dual basis, so integrating its primal basis gives the unique weights.
\writeup (a) Lagrange interpolation gives $P=\sum_iP(x_i)L_i$, hence the weights $w_i=\int_0^1L_i$. Uniqueness follows by testing each $L_i$.

(b) The three polynomials are $2(X-1/2)(X-1)$, $-4X(X-1)$ and $2X(X-1/2)$. Their integrals are $1/6,2/3,1/6$. The rule already integrates degrees zero through two, and at $X^3$ it gives $(0+4/8+1)/6=1/4$, the exact integral. Linearity proves exactness through degree three. Appendix~\ref{linear:ch:inner} discusses nodes chosen for higher exactness.
\end{solution}
\begin{solution}{linear:ex:trig}
\idea Convert zero interpolation data into too many roots of an ordinary polynomial.
\writeup The $2N+1$ exponentials are orthonormal, hence independent; Volume~5, Chapter~6 supplies the circle integral. Multiplying $f$ by $e^{iNx}$ gives $Q(e^{ix})$ with $\deg Q\le2N$. Distinct nodes in $[0,2\pi)$ give distinct circle points. If all evaluation data vanish, $Q$ has $2N+1$ distinct roots, so $Q=0$ and $f=0$. The evaluation map between spaces of dimension $2N+1$ is injective and hence bijective.
\end{solution}
\begin{solution}{linear:ex:lights3}
\idea Single-light targets form a basis; symmetry reduces their exact press certificates to three cases.
\writeup Number cells $(r,c)$ with $1\le r,c\le3$. For a corner target $(1,1)$, press $(1,1),(1,3),(3,1),(2,3),(3,2)$. The corner is toggled once and each other cell twice, as listing its neighbouring pressed cells verifies. For the edge target $(1,2)$, press $(2,2),(3,1),(3,2),(3,3)$: the target is toggled once and all others an even number of times. For the centre target, press $(1,2),(2,1),(2,2),(2,3),(3,2)$; the centre is toggled five times and outer cells twice or zero times. These are exact parity certificates, not approximations. Rotating and reflecting them handles all nine targets. Adding the appropriate press vectors therefore solves any board. The image contains a basis, so the map is invertible and the press pattern is unique modulo double presses.
\end{solution}
\begin{solution}{linear:ex:eventown}
\idea The parity Gram matrix has rank $m$ or $m-1$, leaving little room for more than $n$ vectors.
\writeup With indicators as rows of $M$, the Gram matrix is $G=MM^{\mathsf T}=J-I$, so $\operatorname{rank}G\le n$. In characteristic two, $G^2=mJ+I$. If $m$ is even, $G^2=I$ and $m\le n$. If $m$ is odd, constants lie in its kernel; its leading $(m-1)$ block is invertible by the even case, so its rank is exactly $m-1$ and $m\le n+1$. This also includes $m=1$, where the empty minor has rank zero.
\end{solution}
\begin{solution}{linear:ex:singularspace}
\idea A common kernel vector produces many singular matrices; evaluation at that vector bounds an invertible subspace.
\writeup (a) Matrices with zero first column form a dimension-$n(n-1)$ subspace, all killing $e_1$. (b) On a subspace $L$ whose nonzero matrices are invertible, the linear map $A\mapsto Ae_1$ is injective. Its target has dimension $n$, so $\dim L\le n$. In $M_2(\R)$, the matrices $\begin{pmatrix}a&-b\\b&a\end{pmatrix}$ have determinant $a^2+b^2$, proving the dimension-two example. Volume~3, Chapter~12, gives the quaternionic dimension-four analogue.
\end{solution}
\begin{solution}{linear:ex:hyperplanes}
\idea Factor through the map recording the finitely many given measurements.
\writeup Let $F(v)=(\varphi_i(v))_{i=1}^k$. A linear combination of those forms plainly kills $\ker F$. Conversely, if $\varphi$ kills $\ker F$, it induces a form on $\operatorname{Im}F$ by $F(v)\mapsto\varphi(v)$. Extend that form to $\mathbb K^k$ by a complement; it has coordinates $c_i$, giving $\varphi(v)=\sum_ic_i\varphi_i(v)$. This proof only uses finite dimension in the target, so works for arbitrary $V$ as well. In finite dimension it is exactly the transpose-image annihilator identity.
\end{solution}
\begin{solution}{linear:ex:hamel}
\idea Rational-linearity permits arbitrary basis values, but boundedness would force real-linearity.
\writeup Extend the assignment $f(b_1)=b_1$ and $f(b)=0$ for other basis elements rational-linearly. There are at least two basis elements. The values at $b_1,b_2$ would force a putative real slope to equal both one and zero, so the additive function is not $cx$.

Suppose an additive $f$ is bounded by $M$ on $[u,u+\delta]$, with $\delta>0$. Differences show $|f(t)|\le2M$ for $|t|\le\delta$. Set $c=f(1)$ and $g(x)=f(x)-cx$. Additivity gives rational-linearity, so $g$ vanishes on $\Q$ and is bounded by $B=2M+|c|\delta$ on that small interval. For each $x$ and positive integer $n$, choose rational $r$ with $|x-r|<\delta/n$. Then $|g(x)|=|g(n(x-r))|/n\le B/n$. Let $n\to\infty$, obtaining $g=0$. Thus the constructed function is unbounded on every interval of positive length. Its graph is also dense: the graph points over the rational span of $b_1,b_2$ are the image of dense $\Q^2$ under the invertible real matrix with columns $(b_1,b_1)$ and $(b_2,0)$.
\end{solution}
\begin{solution}{linear:ex:fisher}
\idea The common intersection term and positive diagonal corrections make the real Gram matrix nonsingular.
\writeup If $m\le1$, the conclusion follows from $n\ge1$. Otherwise each $d_i=|C_i|-\lambda$ is non-negative. Two zero values would imply $C_i=C_i\cap C_j=C_j$, contradicting distinctness, so at most one vanishes. For the indicator Gram matrix $G=\lambda J+D$,
\[x^{\mathsf T}Gx=\lambda(\sum_ix_i)^2+\sum_id_ix_i^2.
\]
If this is zero, all coordinates with $d_i>0$ vanish, leaving at most one coordinate; the first term then forces that coordinate to vanish too. Thus $G$ is positive definite and has rank $m$. Since it is $MM^{\mathsf T}$ for an $m\times n$ matrix, $m\le n$. The design-theoretic Fisher inequality has a dual incidence formulation, developed in Volume~3.
\end{solution}
\begin{solution}{ex:v8cAa}
\idea The quotient forgets exactly the polynomials with two zero endpoint values.
\writeup The generators of $W$ are independent and vanish at zero and one. Evaluation $E(P)=(P(0),P(1))$ is onto, using the linear interpolant, so its kernel has dimension two and equals $W$. Therefore $W^\circ=\operatorname{span}(\varepsilon_0,\varepsilon_1)$ and $E$ induces the quotient isomorphism. The representative is $P(0)(1-X)+P(1)X$. For the given cubic the values are one and ten, giving $1+9X$.
\end{solution}
\begin{solution}{ex:v8cAb}
\idea A complement chooses one representative over each point of the fixed complement $C$.
\writeup Projection along $W$ maps any other complement $D$ isomorphically onto $C$. Its inverse has the form $c\mapsto L(c)+c$, with $L$ linear and unique. Conversely every graph of a linear $L$ is a complement, by direct decomposition and zero intersection. A map from dimension $n-k$ to dimension $k$ over $\mathbb F_q$ has $k(n-k)$ freely chosen matrix entries, giving $q^{k(n-k)}$ graphs.
\end{solution}
\begin{solution}{ex:v8cAc}
\idea The third equation repeats the sum of the first two, so their discrepancy is the obstruction.
\writeup The transpose kernel is $\mathbb K(-1,-1,1)$. Thus solvability is equivalent to $d=b_3-b_1-b_2=0$, in which case the unique solution is $(b_1,b_2)$. If $d\ne0$, take $y=d^{-1}(-1,-1,1)$, which satisfies $A^{\mathsf T}y=0$ and $y^{\mathsf T}b=1$. These formulas work in characteristic two as well, with signs interpreted in the field.
\end{solution}
\begin{solution}{ex:v8cAd}
\idea Zero values on the two coordinate axes remove every term except $XY$.
\writeup Write $P=a+bX+cY+dX^2+eXY+fY^2$. Vanishing at $(0,0)$ gives $a=0$. At $(1,0),(2,0)$ it gives $b+d=0$ and $2b+4d=0$, hence $2d=0$ and $b=d=0$. The two analogous $Y$ equations give $c=f=0$. Finally $(1,1)$ gives $e=0$. Thus the evaluation map is injective between six-dimensional spaces and hence bijective. Division by two was essential; in characteristic two the listed zero and two nodes would also coincide.
\end{solution}

\section*{Sources and further reading}

\begin{itemize}
\item H.~Grassmann, \emph{Die lineale Ausdehnungslehre}, Wigand, Leipzig, 1844; English translation by L.~C.~Kannenberg,
\emph{A New Branch of Mathematics}, Open Court, 1995. The story of its reception is told in D.~Fearnley-Sander,
\emph{Hermann Grassmann and the creation of linear algebra}, American Mathematical Monthly 86 (1979), 809--817.
\item G.~Peano, \emph{Calcolo geometrico secondo l'Ausdehnungslehre di H.~Grassmann}, Bocca, Turin, 1888. Chapter~IX gives
the axioms of a vector space.
\item S.~Axler, \emph{Linear Algebra Done Right}, 4th ed., Springer, 2024. Chapter~3 for quotients, duality and the
transpose, in the basis-free spirit of this appendix.
\item P.~R.~Halmos, \emph{Finite-Dimensional Vector Spaces}, 2nd ed., Van Nostrand, 1958. The classic treatment of
annihilators and the double dual.
\item L.~Babai and P.~Frankl, \emph{Linear Algebra Methods in Combinatorics}, preliminary version, University of Chicago,
1992. Chapter~1 opens with Oddtown and Eventown.
\item E.~R.~Berlekamp, \emph{On subsets with intersections of even cardinality}, Canadian Mathematical Bulletin 12 (1969),
471--474.
\item K.~Sutner, \emph{Linear cellular automata and the Garden-of-Eden}, Mathematical Intelligencer 11 (1989), 49--53, for
the linear algebra of Lights Out on general boards.
\item H.~Flanders, \emph{On spaces of linear transformations with bounded rank}, Journal of the London Mathematical Society
37 (1962), 10--16: a subspace of singular $n\times n$ matrices has dimension at most $n(n-1)$, which shows that
Exercise~\ref{linear:ex:singularspace}(a) is optimal.
\item R.~A.~Fisher, \emph{An examination of the different possible solutions of a problem in incomplete blocks}, Annals of
Eugenics 10 (1940), 52--75.
\end{itemize}
